\begin{frame}{Espectro de energia}
  \begin{equacao}
    \[
      E_n = \left(n+\frac12\right)\hbar\omega,
      \qquad n = 0,1,2,\ldots
    \]
  \end{equacao}

  \begin{columns}[T,onlytextwidth]
    \begin{column}{0.53\textwidth}
      \begin{itemize}
        \item Níveis \alert{igualmente espaçados}:
              $\Delta E = \hbar\omega$.
        \item O estado fundamental tem energia \alert{não nula}.
      \end{itemize}
      \begin{destaque}{Energia de ponto zero (via Lua)}
        Para $\hbar\omega = 250\ \mathrm{meV}$:
        \[
          E_0 = \frac12\hbar\omega = \energia{0}{250}
        \]
      \end{destaque}
    \end{column}
    \begin{column}{0.44\textwidth}
      \begin{destaque}[fisicaamber]{Consequência}
        Mesmo a $T=0$ o oscilador não está em repouso. Essa energia
        residual é exigida pelo princípio da incerteza e tem efeitos
        mensuráveis em vibrações moleculares e em qubits.
      \end{destaque}
    \end{column}
  \end{columns}
\end{frame}
